% 第10章: 設計のコツ集
\section{設計のコツ集}

この章では、実装でよくある疑問に短いQ\&A形式で答えます。

%% ── エラー処理 ──────────────────────────────
\begin{qabox}{エラーが起きたらどうする？}

\texttt{init()}は\texttt{bool}を返します。失敗したら\texttt{false}を返しましょう。

\begin{lstlisting}
bool TempSensorModule::init() {
    if (!checkSensor()) {
        Serial.println("[TempSensor] init failed: sensor not found");
        return false;  // 失敗を通知
    }
    return true;
}
\end{lstlisting}

\texttt{updateInput()}/\texttt{updateOutput()}でエラーが起きた場合は、Dataのフラグで通知します。

\begin{lstlisting}
void TempSensorModule::updateInput(SystemData& data) {
    int raw = analogRead(_config.analogPin);
    if (raw < 0) {
        data.temp.isValid = false;  // エラーフラグ
        return;
    }
    data.temp.temperature = raw * _config.conversionFactor;
    data.temp.isValid = true;
}
\end{lstlisting}

\end{qabox}

%% ── バス共有 ──────────────────────────────
\begin{qabox}{SPIやI2Cを複数モジュールで共有するには？}

バス（SPI、I2C）のインスタンスは \texttt{main.cpp} で作り、コンストラクタの引数でモジュールに渡します。

\begin{lstlisting}
// main.cpp
static TwoWire i2cBus = TwoWire(0);

void setup() {
    i2cBus.begin(SDA_PIN, SCL_PIN);  // バス初期化は先に！

    // モジュール初期化は後
    tempSensor.init();
    display.init();
}
\end{lstlisting}

\begin{lstlisting}
// モジュール側 — ポインタで受け取る
class TempSensorModule : public IModule {
private:
    TwoWire* _wire;

public:
    TempSensorModule(const TempSensorConfig& config,
                     TwoWire* wire)
        : _config(config), _wire(wire) {}
};
\end{lstlisting}

\begin{cautionbox}[バス初期化の順番]
\texttt{bus.begin()}は必ず\texttt{setup()}内で、全モジュールの\texttt{init()}より前に呼んでください。モジュールの\texttt{init()}内で\texttt{bus.begin()}を呼ぶと、2つ目のモジュールが初期化するときにバスがリセットされる可能性があります。
\end{cautionbox}

SPI使用時は排他制御を忘れないでください。

\begin{lstlisting}
void updateOutput(SystemData& data) {
    _spi->beginTransaction(SPISettings(1000000, MSBFIRST, SPI_MODE0));
    // SPI通信
    _spi->endTransaction();
}
\end{lstlisting}

\end{qabox}

%% ── init失敗後の復帰 ──────────────────────
\begin{qabox}{init失敗したモジュールは復帰できる？}

できます。ロジックフェーズで定期的にinit()を再試行します。

\begin{lstlisting}
// ロジックフェーズ内
void applyLogic(SystemData& data) {
    // 無効化されたモジュールの再init（10秒ごと）
    static ModuleTimer retryTimer;
    if (retryTimer.getNowTime() >= 10000) {
        retryTimer.setTime();
        for (int i = 0; i < INPUT_COUNT; i++) {
            if (!inputModules[i]->enabled) {
                if (inputModules[i]->init()) {
                    inputModules[i]->enabled = true;
                    Serial.println("[System] module recovered");
                }
            }
        }
    }
}
\end{lstlisting}

接触不良で一時的に通信が切れたセンサーが、再接続後に自動復帰します。

\end{qabox}

%% ── ログ出力 ──────────────────────────────
\begin{qabox}{ログはどう書くのがいい？}

統一フォーマットを使いましょう。

\begin{lstlisting}
// フォーマット: [モジュール名] メッセージ
Serial.println("[TempSensor] init OK");
Serial.println("[TempSensor] read failed: timeout");
Serial.println("[LED] state changed: ON");
\end{lstlisting}

角括弧でモジュール名を付けると、シリアルモニターでどのモジュールの出力か一目で分かります。

\end{qabox}

%% ── Config変更 ──────────────────────────────
\begin{qabox}{実行中にConfigを変更したいときは？}

Config変更はロジックフェーズでのみ行います。

\begin{badexample}{入力フェーズでConfig変更}
\begin{lstlisting}
// BLEモジュールのupdateInput()内で直接変更 → 危険！
void BleModule::updateInput(SystemData& data) {
    if (newConfigReceived) {
        ledModule._config.pin = newPin;  // 他モジュールを直接操作
    }
}
\end{lstlisting}
\end{badexample}

\begin{goodexample}{SystemData経由でロジックフェーズに渡す}
\begin{lstlisting}
// BLEモジュール — 受信データをSystemDataに書くだけ
void BleModule::updateInput(SystemData& data) {
    if (newConfigReceived) {
        data.ble.requestedLedPin = newPin;
        data.ble.configUpdateRequested = true;
    }
}

// ロジックフェーズ — ここでConfig変更を実行
void applyLogic(SystemData& data) {
    if (data.ble.configUpdateRequested) {
        // 安全にConfig変更を適用
        data.ble.configUpdateRequested = false;
    }
}
\end{lstlisting}
\end{goodexample}

\end{qabox}

%% ── ビルドエラー ──────────────────────────
\begin{qabox}{ビルドが通らない！よくあるエラーは？}

初めてモジュールを追加したとき、コンパイルエラーに悩むことが多いです。よくあるエラーとその対処法をまとめます。

\subsection{「No such file or directory」}

\begin{lstlisting}[language={},numbers=none,backgroundcolor=\color{warnred!5}]
fatal error: SystemData.h: No such file or directory
\end{lstlisting}

\texttt{lib/} 内のコードから \texttt{include/} のヘッダーが見つかっていません。\texttt{platformio.ini} の \texttt{build\_flags} に \texttt{-I include} が入っているか確認してください。

\begin{lstlisting}[language={},keywordstyle={},commentstyle=\color{gray}\itshape,morecomment={[l]{;}},numbers=none]
build_flags =
    -I include            ; これがないとlib/からinclude/が見えない
\end{lstlisting}

\subsection{「does not name a type」}

\begin{lstlisting}[language={},numbers=none,backgroundcolor=\color{warnred!5}]
error: 'SystemData' does not name a type
\end{lstlisting}

ヘッダーファイルで \texttt{struct SystemData;} の前方宣言を忘れている、または \texttt{.h} で \texttt{SystemData.h} をインクルードしてしまい循環依存が発生しています。以下を確認してください。

\begin{itemize}
  \item \texttt{XxxModule.h} — \texttt{struct SystemData;} の前方宣言があるか
  \item \texttt{XxxModule.h} — \texttt{SystemData.h} をインクルードしていないか
  \item \texttt{XxxModule.cpp} — \texttt{\#include "SystemData.h"} があるか
\end{itemize}

\subsection{「is not a member of 'SystemData'」}

\begin{lstlisting}[language={},numbers=none,backgroundcolor=\color{warnred!5}]
error: 'struct SystemData' has no member named 'temp'
\end{lstlisting}

新しいモジュールのDataを \texttt{SystemData.h} に追加し忘れています。\texttt{include/SystemData.h} を開いて、自分のモジュールの \texttt{\#include} と \texttt{Data} メンバを追加してください（実践の章ステップ5参照）。

\subsection{「marked 'override', but does not override」}

\begin{lstlisting}[language={},numbers=none,backgroundcolor=\color{warnred!5}]
error: 'bool XxxModule::init()' marked 'override', but does not
override any member functions
\end{lstlisting}

\texttt{IModule} を継承していないか、メソッドのシグネチャ（引数・戻り値）が合っていません。\texttt{: public IModule} の継承と、\texttt{init()}/\texttt{updateInput()}/\texttt{updateOutput()} の引数が正しいか確認してください。

\begin{notebox}[エラーメッセージの読み方]
コンパイルエラーは上から順に読みます。最初のエラーが原因で後続のエラーが連鎖的に発生することが多いため、まず最初の1つを修正してから再ビルドしてください。
\end{notebox}

\end{qabox}

%% ── deinit ──────────────────────────────
\begin{qabox}{deinit()はいつ使うの？}

\texttt{deinit()}は、モジュールが確保したリソースを解放するためのメソッドです。IModuleにデフォルトの空実装があるので、必要なモジュールだけオーバーライドします。

\begin{lstlisting}
class CameraModule : public IModule {
public:
    void deinit() override {
        // フレームバッファを解放
        if (_fb) {
            esp_camera_fb_return(_fb);
            _fb = nullptr;
        }
        esp_camera_deinit();
        Serial.println("[Camera] deinit done");
    }
};
\end{lstlisting}

主な用途は以下の通りです。

\begin{itemize}
  \item ディープスリープに入る前のリソース解放
  \item モジュールの動的な停止・再起動
  \item テスト時の後片付け
\end{itemize}

\end{qabox}
